\title{Controlling Text-to-Image Diffusion by Orthogonal Finetuning}
When considering an $r$-block diagonal structure, the number of free parameters becomes $\thickmuskip=2mu \medmuskip=2mu d(d/r-1)/2$, reducing the complexity to $\thickmuskip=2mu \medmuskip=2mu \mathcal{O}(d^2/r)$. \section{Intriguing Insights and Discussions}

\textbf{OFT’s architectural flexibility}. \end{itemize}

At their core, both tasks focus on the challenge of finetuning text-to-image diffusion models in a manner that maintains the original generative capabilities acquired during pretraining. Expanded results are reported in Appendix~\ref{app:more_qualitative},~\ref{app:more_control_task},~\ref{app:failure}. A further reduction is possible by sharing the block matrices across all positions, that is, enforcing $\thickmuskip=2mu \medmuskip=2mu\bm{R}_i=\bm{R}_j$ for any $i\neq j$, bringing parameter count down to $\mathcal{O}(d^2/r^2)$. Beyond dense (fully-connected) layers, OFT extends naturally to convolutional layers, with the block-diagonal design of $\bm{R}$ taking on meaningful interpretations in the convolutional context—contrasting with LoRA's limitations. While conventional finetuning approaches offer maximum adaptability, they tend to be unstable and are susceptible to model collapse. Furthermore, we show that OFT generalizes effectively to additional control settings—including dense pose-to-human body, sketch-to-image, and depth-to-image tasks—as detailed in Appendix~\ref{app:more_control_task}. As indicated in Table~\ref{table:dreambooth_final}, both COFT and OFT substantially surpass DreamBooth and LoRA with respect to the DINO and CLIP-I metrics, and also deliver marginally superior or at least equivalent results for text fidelity and diversity. Concretely, we define $\bm{R}$ by the transformation $\thickmuskip=2mu \medmuskip=2mu \bm{R}=(\bm{I}+\bm{Q})(I-\bm{Q})^{-1}$, where the matrix $\bm{Q}$ is skew-symmetric and fulfills $\thickmuskip=2mu \medmuskip=2mu \bm{Q}=-\bm{Q}^\top$. \item Our investigation applies OFT to two scenarios: subject-driven image synthesis and controlled generation. \textbf{Why OFT converges faster}. To further bolster the robustness of finetuning, we introduce Constrained Orthogonal Finetuning (COFT), which applies a supplementary radius constraint to the hypersphere. Comparisons to LoRA are controlled by limiting OFT and COFT modifications to the same network layers utilized by LoRA. CLIP-T computes the mean cosine similarity between CLIP embeddings of text prompts and corresponding generated images. Remarkable advancements~\cite{saharia2022photorealistic,ramesh2022hierarchical,rombach2022high,gu2022vector,nichol2022glide} have been achieved in the realm of text-to-image synthesis, which can be attributed in large part to developments in diffusion-based generative modeling~\cite{ho2020denoising,song2021score,song2021denoising,dhariwal2021diffusion} as well as in vision-language joint representations~\cite{su2020vlbert,lu2019vilbert,radford2021learning,wang2021simvlm,li2022blip,li2023blip,alayrac2022flamingo,schuhmann2022laion}. Further inspiration is drawn from orthogonal over-parameterized training~\cite{liu2021orthogonal}, a method in which classification networks are trained from random initialization via orthogonal transformations. By contrast, ControlNet demonstrates abrupt controllability that manifests after epoch 8, aligning with the sudden convergence effect noted by \cite{zhang2023adding}. Notably, OFT also demonstrates heightened data efficiency, reaching convergence with considerably fewer training examples and epochs. Crucially, OFT does not add extra inference computation, thus ensuring models remain efficient for deployment. Human evaluation results in Appendix~\ref{app:human} further substantiate OFT’s superiority. The method is specifically applied to two scenarios: subject-driven generation and controllable generation. In the case of GLIDE~\cite{nichol2022glide} and Imagen~\cite{saharia2022photorealistic}, diffusion models are trained directly in the pixel domain; GLIDE leverages paired text-image data to jointly optimize the text encoder and the diffusion prior, whereas Imagen utilizes a pretrained, fixed text encoder. The corresponding results can be seen in Figure~\ref{fig:teaser} and Figure~\ref{fig:db_convergence}. Second, OFT holds promise for compositionality: the orthogonal matrices yielded from various OFT finetuning tasks can be successively multiplied, with the product retaining orthogonality. The OFT framework is depicted in Figure~\ref{fig:block}. \textbf{Subject-driven generation}. LoRA utilizes a low-rank matrix addition to the pretrained weights, which limits drastic modifications to the pretrained weight matrix. COFT effectively unifies two direct constraints: a restriction on hyperspherical energy and a bounded departure from the original pretrained parameters. Next, we compare the parameter efficiency of OFT against LoRA. To address this, we propose maintaining the angles among neurons, operating under the premise that these angular relationships are central to the neural network’s knowledge representation. \textbf{Why not minimize hyperspherical energy}. The visualizations in Figure~\ref{fig:dreambooth_final} indicate that both OFT and COFT excel in preserving semantic details of the subject, whereas LoRA frequently struggles to maintain subject identity (\emph{e.g.}, LoRA entirely fails in the bowl case). We investigate three demanding controllable generation scenarios in the main text: converting Canny edges to images~(C2I) using the COCO dataset~\cite{lin2014microsoft}, translating segmentation maps to images~(S2I) on the ADE20K dataset~\cite{zhou2017scene}, and generating facial images from landmarks~(L2F) using the CelebA-HQ dataset~\cite{karras2018progressive,xia2021tedigan}. We further substantiate this claim with an experiment that assesses the regularization benefit of enforcing orthogonality. Notably, LoRA and OFT illustrate fundamentally different mechanisms for efficient parameterization: LoRA leverages low-rank structure for parameter reduction, whereas OFT capitalizes on structural sparsity (\emph{i.e.}, block-diagonal orthogonal matrices). For DreamBooth and LoRA, experimental procedures and hyperparameters are adopted from their respective primary sources, using recommended optimal values. For S2I, metrics include mean IoU as well as both mean and overall accuracy. CLIP-I reports the mean pairwise cosine similarity between the CLIP embeddings of real and generated images. However, empirical results suggest that this constraint does not adversely impact model performance. Employing a controllable image generation protocol in the style of ControlNet~\cite{zhang2023adding}, we present results in Figure~\ref{fig:ortho}. To better illustrate the advantages of OFT, we present a selection of randomly chosen subject-driven generation samples in Figure~\ref{fig:dreambooth_final}. While this strategy enforces an intuitively reasonable restriction, it may nonetheless permit more deviation than is desirable when finetuning large-scale models. \textbf{Qualitative comparison}. Despite its advantages, OFT brings forth several intriguing open questions. This metric estimates the control cue from the generated output and compares it with the initial control input, offering an effective means of validating OFT’s superior controllability characteristics. These observations demonstrate that our OFT framework combines rapid convergence with high stability in subject-driven generative tasks. The visual results in Figure~\ref{fig:face_convergence_qual} indicate that both OFT and COFT exhibit stable convergence, progressively aligning the synthesized face pose to the specified control landmarks. \textbf{Connection to LoRA}. For LoRA, if the low-rank parameter is $r'$, the total trainable parameter count is given by $\thickmuskip=2mu \medmuskip=2mu r'(d+n)$. In the C2I setting, OFT and COFT effectively generate images that semantically align with rough Canny edges, outperforming T2I-Adapter and LoRA, which show inferior results. To this end, COFT adopts the following inference step:

\begin{equation}\label{eq:coft_general}
    \footnotesize
    \bm{z}=\bm{W}^\top\bm{x}=(\bm{R}\cdot\bm{W}^0)^\top\bm{x},~~~\text{s.t.}~\bm{R}^\top\bm{R}=\bm{R}\bm{R}^\top=\bm{I},~\norm{\bm{R}-\bm{I}}\leq \epsilon
\end{equation}

where the orthogonality of $\bm{R}$ is enforced alongside a bound on its distance from the identity matrix by at most $\epsilon$. Importantly, this reduced sensitivity to the finetuning iteration count relieves the burden of hyperparameter tuning for individual subjects; both OFT and COFT allow for the direct specification of a relatively high number of iterations without the need for fine-tuning. Increasing $\epsilon$ allows greater adaptation and causes the COFT model to more closely approximate the OFT-finetuned solution. Parallel to LoRA's rank hyperparameter $r'$, OFT introduces a diagonal block parameter $r$, which also curtails the number of tunable parameters. This underscores that modifying the directions of neurons is fundamental for semantic alterations in generated images—a principle at the heart of both subject-driven and controllable generation. For enhanced computational efficiency, our method employs Cayley parameterization to construct orthogonal matrices. \subsection{Subject-driven Generation}

\textbf{Configuration}. Achieving a balance between adaptability and stability is pivotal in finetuning, and we contend that our OFT framework adeptly navigates this trade-off. Second, OFT introduces an explicit constraint to limit deviation from the pretrained model, giving rise to a constrained variant—an element that is absent from \cite{liu2021orthogonal}. Although OFT achieves strong results, our observations indicate that COFT further enhances finetuning stability owing to its direct constraint on parameter deviation. From these results, it is evident that both COFT and OFT achieve convergence markedly faster than competing approaches. In comparison to prior approaches, OFT offers enhanced control and greater stability during finetuning, while requiring fewer tunable parameters. Standard approaches to finetuning either make incremental adjustments to neuron weights using a modest learning rate (\emph{e.g.}, \cite{ruiz2023dreambooth}), or introduce lightweight modules through re-parameterized neuron weights (\emph{e.g.}, \cite{hulora2022,zhang2023adding}). The superior data-efficiency of OFT is further illustrated in Figure~\ref{fig:teaser}, where the right-side example shows OFT achieving satisfactory convergence using only 5\% of the ADE20K dataset, outperforming ControlNet and LoRA under limited data conditions. DreamBooth~\cite{ruiz2023dreambooth} and LoRA~\cite{hulora2022} are chosen as baseline methods. For classification tasks, having non-redundant neurons is advantageous. During training, the standard inner product computes the feature map (where $z$ represents the convolution output for kernel $\bm{w}$ and sliding window input $\bm{x}$). \item To quantitatively assess controllable generation, we propose a novel metric called control consistency. To further highlight the significance of neuronal angular components, we run a simplified experiment illustrated in Figure~\ref{fig:toy_ae}, where a typical convolutional autoencoder is trained on a collection of flower images. Ensuring the orthogonality of $\bm{R}$ can be accomplished through differentiable orthogonalization approaches, as described in~\cite{liu2021orthogonal,lezcano2019cheap}. Through an extensive series of empirical analyses, we show that OFT delivers notable gains over established methods, excelling in generation fidelity, training speed, and finetuning robustness. The Cayley parameterization does restrict generated matrices to the special orthogonal group, specifically orthogonal matrices with determinant $1$. First, the technique enforces orthogonality via Cayley parametrization, which depends on matrix inversion—posing some constraints on scalability. In contrast to ControlNet, the OFT approach incurs no extra computational cost during inference for the finetuned model. For tasks centered on subject-driven image synthesis, we follow the methodology outlined in DreamBooth~\cite{ruiz2023dreambooth}. To fulfill this requirement, we set the initial value of the orthogonal matrix $\bm{R}$ as the identity matrix, ensuring that the process begins from the original pretrained weights. In contrast, OFT introduces a distinct type of restriction, enforcing invariance of pairwise neuron similarity by demanding that $\thickmuskip=2mu \medmuskip=2mu \| \text{HE}(\bm{M})-\text{HE}(\bm{M}^0) \|=0$, where $\text{HE}(\cdot)$ computes the hyperspherical energy of a given weight matrix. Concretely, OFT seeks to optimize the orthogonal transformation

\begin{equation}\label{eq:oft_general}
    \footnotesize
    \bm{z}=\bm{W}^\top\bm{x}=(\bm{R}\cdot\bm{W}^0)^\top\bm{x},~~~\text{s.t.}~\bm{R}^\top\bm{R}=\bm{R}\bm{R}^\top=\bm{I}
\end{equation}

In this equation, $\bm{W}$ represents the weight matrix adjusted via OFT-finetuning, and $\bm{I}$ stands for the identity matrix. While a straightforward strategy would be to incorporate a regularization term to stabilize hyperspherical energy during finetuning, this approach cannot ensure the actual minimization of hyperspherical energy discrepancy. OFT is formulated as a model-agnostic method, making it applicable to any text-to-image diffusion architecture. To mitigate this limitation, our work focuses on two distinct categories of text-to-image generation tasks:

\begin{itemize}[leftmargin=*,nosep]

    \item \textbf{Subject-driven generation}~\cite{ruiz2023dreambooth}: With access to only a few example images of a particular subject, this task seeks to create novel images of that subject in varied contexts, using accompanying textual guidance. Inspired by this invariance, we introduce Orthogonal Finetuning~(OFT), an approach designed to adapt large text-to-image diffusion models for new tasks while preserving their original hyperspherical energy. Various inversion methods~\cite{choi2021ilvr,dhariwal2021diffusion,ramesh2022hierarchical,gal2022image} allow the surrounding context to be altered while maintaining the integrity of the subject, and \cite{hertz2022prompt} achieves both global and localized edits independently of input masks. This rationale also motivates the prevalent use of small learning rates in practice. The forward computation for OFT can be reformulated as $\thickmuskip=2mu \medmuskip=2mu \bm{z}=(\bm{R}\bm{W}^0)^\top\bm{x}=(\bm{W}^0+(\bm{R}-\bm{I})\bm{W}^0)^\top\bm{x}$, where $\thickmuskip=2mu \medmuskip=2mu (\bm{R}-\bm{I})\bm{W}^0$ serves a similar function to the low-rank modification employed by LoRA. In contrast, our approach is centered on the finetuning of text-to-image diffusion models, with an emphasis on achieving enhanced controllability and superior generative capabilities for downstream tasks. Specifically, for the C2I task, Intersection over Union (IoU) and F1 score are calculated. \item \textbf{Controllable generation}~\cite{zhang2023adding,mou2023t2i}: Here, the model receives an extra control input (such as canny edges or segmentation masks), and is tasked with generating images conforming to both this control cue and a provided text prompt. Experimental results demonstrate that the OFT method surpasses established baselines in both output quality and speed of convergence. For a full orthogonal matrix of dimension $\thickmuskip=2mu \medmuskip=2mu d\times d$, there are $\thickmuskip=2mu \medmuskip=2mu d(d-1)/2$ independent parameters, which yields a computational complexity of $\mathcal{O}(d^2)$. Building upon empirical findings that hyperspherical similarity effectively encapsulates semantic relationships~\cite{liu2017deep,liu2018decoupled,chen2020angular}, we employ hyperspherical energy~\cite{liu2018learning} to analyze the pairwise interactions among neurons within a layer. In the evaluation phase, we reconstruct the feature map using three distinct strategies: (a) the original inner product, (b) isolating magnitude, and (c) using only angular information. Our baseline comparisons include ControlNet~\cite{zhang2023adding}, T2I-Adapter~\cite{mou2023t2i}, and LoRA~\cite{hulora2022}. A randomly initialized network is theoretically associated with low hyperspherical energy on average, and the objective in \cite{liu2021orthogonal} is to sustain this low energy throughout training, as lower hyperspherical energy is linked to improved generalization in classification tasks~\cite{liu2018learning,lin2020regularizing}. It is evident that COFT and OFT facilitate notably stable finetuning of the diffusion model. The block-diagonal form serves as a form of enhanced regularization on the orthogonal matrix. Image-to-image translation is essentially an instance of controllable generation, with many established approaches relying on conditional generative adversarial networks~\cite{isola2017image,zhu2017toward,wang2018high,park2019semantic,choi2018stargan}. DINO uses a similar approach, utilizing ViT S/16 DINO embeddings. In our qualitative comparisons, we concentrate on OFT, COFT, and ControlNet, given that ControlNet most closely approaches our landmark error. By contrast, when the orthogonality constraint is omitted from OFT, image generation fails to achieve realism and loses control efficacy entirely. Analogous to the strategy of zero initialization found in LoRA, it is necessary to guarantee that OFT finetunes the model starting precisely from the pretrained parameter setting $\bm{W}^0$. \subsection{Constrained Orthogonal Finetuning}

A further restriction of OFT's flexibility can be imposed by confining the finetuned model to lie within a constrained vicinity of the pretrained solution. Notably, LoRA requires 20 epochs to match the accuracy OFT attains by epoch 8. As our quantitative measure, we calculate the mean $\ell_2$ distance between the control face landmarks and those predicted by each model. Though the use of block diagonal parametrization partially mitigates this issue, discovering methods to efficiently perform matrix inversion in a differentiable fashion is still unresolved. Instance-specific image variation is enabled in \cite{casanova2021instance}, albeit sometimes at the cost of losing instance-defining nuances. Despite these reductions, the resulting transformation $\bm{R}$ retains its orthogonal nature, ensuring the preservation of hyperspherical energy properties. \item In contrast to traditional finetuning techniques, OFT allows for adjustment of the model while rigorously maintaining the hyperspherical energy. Regarding the L2F task, our approach demonstrates superior precision in pose control for generated facial images, even in scenarios with difficult head orientations. For scenarios requiring controllable generation, we adhere to protocols established by ControlNet~\cite{zhang2023adding} and T2I-Adapter~\cite{mou2023t2i}. To assess convergence rates, we analyze the performance of ControlNet, T2I-Adapter, LoRA, and COFT on the L2F task through both quantitative metrics and qualitative observations. The key mechanism involves learning an orthogonal transformation shared across layers to guarantee invariance in pairwise neuron angles. We highlight that this quantitative method for evaluating controllable generation constitutes a key contribution of our work, enabling accurate assessment of fine-tuned models’ control capabilities, bounded by the accuracy of existing segmentation or detection models. SphereNet~\cite{liu2017deep} provides evidence that enforcing normalization of all network activations onto a unit hypersphere results in both competitive capacity and enhanced generalization, suggesting that neuron directionality is sufficient to encode essential information from data. The outcomes in Figure~\ref{fig:toy_ae} reveal that reconstructing with only the angular information nearly flawlessly recovers the original images, while the magnitude alone fails to convey any significant content. Notably, adapter-based models (such as T2I-Adapter and ControlNet) exhibit slower convergence and suboptimal final performance. Conversely, our OFT approach demonstrates an upward trend in performance as trainable parameters increase, suggesting its potential has not yet been fully realized. The findings in \cite{liu2021orthogonal} indicate that orthogonal transformations possess sufficient expressive capacity to support the training of neural networks that generalize well in classification scenarios. This consistent convergence trajectory enhances OFT’s practicality, as its smoother and more predictable training progression simplifies the process of tuning its hyperparameters. In the L2F setting, we report the average $\ell_2$ distance between the control landmarks and those predicted. Our principal contributions are outlined as follows:

\begin{itemize}[leftmargin=*,nosep]

    \item We introduce Orthogonal Finetuning, a new approach for refining text-to-image diffusion models to enhance their controllability. This conceptual framework underpins OFT: model finetuning need only seek optimal, layer-wise shared orthogonal matrices to reorient neurons within each layer. \subsection{Re-scaled Orthogonal Finetuning}

We introduce a straightforward augmentation of OFT wherein we additionally optimize a scaling factor for every neuron. Thus, employing a more structural metric to capture differences between the finetuned and pretrained models could significantly enhance both the maintenance of original performance and the stability of the finetuning process. Further experimental details and results are provided in Appendix~\ref{app:settings}. To ensure fairness, outputs from each approach are selected randomly. This procedure guarantees inference latency equivalent to that of the original pretrained network. Notably, cosine activation (c) is not employed during training—only the inner product is used. Additionally, from a geometric perspective, OFT is interpretable as tuning parameters on a smooth hypersphere manifold, which tends to offer more favorable optimization characteristics. Specifically, if the number of blocks matches the number of input channels, each block acts on a single neuron channel, resembling the operation of depth-wise convolutions~\cite{chollet2017xception}. It is important to highlight that OFT and COFT maintain a comparable count of trainable parameters to LoRA—substantially fewer than ControlNet—yet they demonstrate greater efficiency during convergence. Finding strategies to boost parameter efficiency with less bias and greater effectiveness stands out as a significant challenge for future work. Alternatively, with shared block parameters across all channels, the orthogonal transformation by $\bm{R}$ becomes a shared operation. Parameter-efficient finetuning strategies, such as adaptation methods (\emph{e.g.}, \cite{hulora2022,houlsby2019parameter,pfeiffer2020adapterfusion}), have been extensively explored within the field of natural language processing. We begin by assessing the stability during finetuning and the convergence rate for DreamBooth, LoRA, OFT, and COFT. Figure~\ref{fig:face_convergence} presents the landmark error trajectories for models fine-tuned across varying numbers of epochs. Our experiments are conducted with Stable Diffusion v1.5~\cite{rombach2022high} as the base pretrained model for text-to-image tasks. As presented in Table~\ref{table:control_gen}, OFT and COFT achieve markedly superior and more precise control relative to alternative approaches. To address the first question, we take guidance from empirical findings in \cite{liu2018decoupled,chen2020angular}, which note that differences in feature angles effectively encapsulate the semantic disparity. Since $\bm{W}^0$ is usually full-rank, OFT reduces to a low-rank update if $\thickmuskip=2mu \medmuskip=2mu \bm{R}-\bm{I}$ is low-rank. We conduct a qualitative assessment of OFT and COFT in relation to leading methods such as ControlNet, T2I-Adapter, and LoRA. Nevertheless, unconstrained alteration of neuron directions risks undermining the pretrained model’s generative efficacy. Incorporating re-scaling into OFT allows for more adaptive finetuning with only minimal extra parameters. DreamBooth~\cite{ruiz2023dreambooth}, by introducing a special token and leveraging a handful of subject images, finetunes the diffusion model with a combination of reconstruction and class-specific prior preservation losses. We will later elaborate on methods for enforcing this constraint efficiently in computational terms. Furthermore, OFT and COFT both enable precise manipulation using text prompts; in contrast, while DreamBooth maintains subject identity, it often fails to condition images on the textual description (\emph{e.g.}, as seen in the first row for the bowl). Across all three control tasks, OFT and COFT consistently yield images of higher qualitative value compared to other advanced baselines, evidencing the practical utility of our OFT framework in controlled image generation. For most experiments we retain the original OFT formulation to isolate the benefits of orthogonal transformations, but results show that this re-scaled variant often performs even better (refer to Appendix~\ref{app:rescaled}). \textbf{Convergence}. Notably, OFT does not increase inference-time computational requirements. Concerning the second question, the most direct method for adjusting neuron directions involves applying simultaneous rotations or reflections to all neurons within a layer—an approach that is inherently equivalent to orthogonal transformation. By 2000 iterations, DreamBooth tends to produce degenerate images, and LoRA no longer generates a yellow shirt—instead rendering yellow fur. Pivotal Tuning~\cite{roich2022pivotal} refines the generator close to an initial inverted latent vector, with regularization to secure identity preservation, and similarly, \cite{nitzan2022mystyle} develops a custom generative face model from personal image collections. Additionally, we design a constrained variant that restricts the angular change from the original pretrained parameters to promote improved stability during finetuning. \subsection{General Framework}

In standard finetuning approaches, models are commonly updated—either wholly or partially—via gradient descent employing a modest learning rate. The output $\thickmuskip=2mu \medmuskip=2mu \bm{z}\in\mathbb{R}^n$ generated by the layer is given as $\thickmuskip=2mu \medmuskip=2mu \bm{z}=\bm{W}^\top\bm{x}$, with the input $\thickmuskip=2mu \medmuskip=2mu \bm{x}\in\mathbb{R}^d$. \section{Experiments and Results}

\textbf{Experimental setup}. However, while this assumption may sometimes prove valid, we contend that merely minimizing the Euclidean distance to the pretrained model is insufficient for thoroughly assessing semantic preservation. To explore this, we consider the Neumann series as an approximation of $\thickmuskip=2mu \medmuskip=2mu \bm{R}=(\bm{I}+\bm{Q})(I-\bm{Q})^{-1}$, yielding $\th

series converges under the operator norm), allowing us to incorporate the constraint $\thickmuskip=2mu \medmuskip=2mu\|\bm{R}-\bm{I}\|\leq\epsilon$ directly within the Cayley transform. Building on the experimental protocols outlined in \cite{zhang2023adding,mou2023t2i}, we employ OFT to finetune pretrained diffusion models, which consistently results in improved controllability, requiring less training data and fewer parameters to adjust. Notably, OFT and COFT are capable of producing realistic, finely detailed images with plausible geometric structures, while utilizing significantly fewer model parameters than ControlNet. Projected gradient descent provides an efficient mechanism to enforce this new constraint on the matrix $\bm{Q}$. Concretely, our modified forward propagation becomes $\thickmuskip=2mu \medmuskip=2mu\bm{z}=(\bm{R}\bm{W}^0\bm{D})^\top\bm{x}$\footnote{\textbf{Errata}: The NeurIPS camera-ready manuscript incorrectly described the re-scaled OFT forward pass as $\thickmuskip=2mu \medmuskip=2mu\bm{z}=(\bm{D}\bm{R}\bm{W}^0)^\top\bm{x}$. With this perspective, it becomes compelling to apply a layer-wise orthogonal transformation to neurons, thereby ensuring the preservation of hyperspherical energy across layers. \textbf{Quantitative comparison}. Two primary distinctions separate OFT from the methodology of \cite{liu2021orthogonal}. Our conjecture is that an optimal finetuned model should maintain minimal deviation in hyperspherical energy relative to its pretrained counterpart. Likewise, T2I-Adapter~\cite{mou2023t2i} focuses on enhancing the controllability of diffusion-generated images by finetuning pretrained diffusion models. Strategies aiming to preserve the original subject include approaches such as~\cite{avrahami2022blended,nichol2022glide}, which make use of user-provided masks as supplemental input. For consistency, all examples derive from the same finetuned model across different approaches; individual text prompts are not separately optimized for final outputs. To address this, we present Orthogonal Finetuning~(OFT), a systematic approach for refining text-to-image diffusion models for specific downstream tasks. Each method was evaluated by repeatedly finetuning the same pretrained model with 30 different random seeds to minimize the impact of randomness. By contrast, OFT and COFT continue to exhibit reliable and robust control over image generation. Additionally, OFT is robust to the number of optimization iterations, maintaining strong performance even with many iterations, whereas DreamBooth and LoRA do not share this trait. All models utilize the same loss criterion as DreamBooth. We observe that conserving this attribute is essential to uphold the semantic generative competence of text-to-image diffusion models. As illustrated in Figure~\ref{fig:toy_ae}, the most direct and impactful means to alter a network's semantics involves adjusting the directions of neurons—a capability that OFT directly targets. In contrast, OFT achieves preservation of pretraining knowledge by enforcing the transformation on the pretrained weight matrix to be orthogonal and full-rank. Empirical evidence in \cite{liu2021orthogonal} supports this improved landscape. To mitigate this inefficiency, we opt to express $\bm{R}$ as a block-diagonal matrix containing $r$ blocks, yielding the form:

\begin{equation}
    \footnotesize
    \bm{R}=\text{diag}(\bm{R}_1,\bm{R}_2,\cdots,\bm{R}_r)=\begin{bmatrix}
    \bm{R}_1\in \text{O}(\frac{d}{r}) &  &\\
    &\ddots & \\
    & & \bm{R}_r\in \text{O}(\frac{d}{r})
    \end{bmatrix}\in \text{O}(d)
\end{equation}

Here, $\text{O}(d)$ refers to the group of $d$-dimensional orthogonal matrices, with $\thickmuskip=2mu \medmuskip=2mu \bm{R}\in\mathbb{R}^{d\times d}$ and each block $\thickmuskip=2mu \medmuskip=2mu \bm{R}_i\in\mathbb{R}^{d/r\times d/r}$ for all $i$, all being orthogonal. Each method finetunes Stable Diffusion~(SD) v1.5 on the respective datasets for 20 epochs. To impose a stricter control, LoRA supplements the optimization with an explicit low-rank constraint on the update: $\thickmuskip=2mu \medmuskip=2mu \text{rank}(\bm{M} - \bm{M}^0)=r'$, for a pre-specified, small $r'$. Additionally, \cite{liu2021orthogonal} demonstrates the sufficiency of orthogonal transformations in the context of neural network learning. If $r$ and $r'$ scale with $d$—for example, if $\thickmuskip=2mu \medmuskip=2mu r=r'=\alpha d$ where $\thickmuskip=2mu \medmuskip=2mu 0<\alpha\leq 1$ is a constant—then LoRA exhibits parameter complexity $\mathcal{O}(d^2+dn)$, whereas OFT's parameter count scales as $\mathcal{O}(d)$. Our findings reveal an inherent compromise between flexibility and regularity when finetuning models. While LoRA surpasses ControlNet in S2I, its results lag in C2I and L2F, and even with meticulous tuning, LoRA tends to reach a lower overall performance threshold. Unlike the original OFT forward computation in Eq.~\eqref{eq:oft_general}, where only $\bm{R}$ was updated during training, both $\bm{D}$ and $\bm{R}$ are now optimized. Appendix~\ref{app:settings} provides a comprehensive explanation of the consistency metrics. For the S2I task, LoRA fails to synthesize images conforming to the provided segmentation map, whereas ControlNet, OFT, and COFT successfully regulate image generation as intended. Furthermore, for COFT, given a suitable $\epsilon$, both the learning rate and the number of iterations become largely non-critical, simplifying parameter selection. To ensure that the orthogonal matrix $\bm{R}$ is initialized as the identity, we simply set $\bm{Q}$ initially to be all zeros. \end{document} This process can be interpreted as modifying the standard basis for neurons within each layer. The ControlNet framework~\cite{zhang2023adding} improves the controllability of pretrained diffusion models through finetuning and the inclusion of additional conditioning inputs, demonstrating significant advances in controllable generation. OFT operates within the DreamBooth paradigm but, rather than straightforward parameter updating with a

\textbf{Controllable generation}. Our findings indicate that the OFT framework consistently delivers improved stability, faster convergence, and superior final performance compared to previous approaches. First, while \cite{liu2021orthogonal} is oriented toward minimizing hyperspherical energy, OFT is specifically designed to conserve the pretrained model’s hyperspherical energy, thereby protecting the pretrained structural properties during finetuning; minimizing hyperspherical energy during diffusion model finetuning could otherwise compromise established semantic structures. A key open question is how to efficiently steer or modulate these sophisticated models for various downstream applications. The OFT approach is, in theory, applicable to a wide array of neural network structures. Apart from diffusion strategies, research into text-to-image synthesis has included generative adversarial networks~\cite{reed2016generative,zhang2017stackgan,xu2018attngan,li2019controllable} and autoregressive models~\cite{ramesh2021zero,ding2021cogview,wu2022nuwa,yuscaling}. Supplemental experimental outcomes are available in Appendix~\ref{app:settings},~\ref{app:coft_oft},~\ref{app:more_qualitative},~\ref{app:failure}. Specifically, hyperspherical energy aggregates the hyperspherical similarity values (such as cosine similarity) between all pairs of neurons in a given layer, thus reflecting how uniformly the neurons are distributed over the unit hypersphere~\cite{liu2021learning}. A configuration that attains minimum hyperspherical energy places neurons evenly distributed on the hypersphere, which is not suitable for finetuning as it could compromise the integrity of the information learned during pretraining. We also assess LPIPS, reporting the mean cosine similarity across pairs of generated images with the same subject and prompt. Although one could contemplate more nuanced angular transformations that independently rotate each neuron, our analysis suggests that orthogonal transformations offer an optimal trade-off between adaptability and structural regularity. Lastly, the efficiency of OFT, in terms of parameters, is closely tied to the block diagonal design, which unavoidably introduces certain biases and can restrict flexibility. As illustrated in Figure~\ref{fig:eps_coft}, varying $\epsilon$ significantly influences COFT's effectiveness. \textbf{Quantitative comparison}. While the latter constraint is usually present only in an implicit form within existing finetuning techniques, COFT makes this deviation explicit. This translates into the equivalent restriction $\thickmuskip=2mu \medmuskip=2mu \|\bm{Q}-\bm{0}\|\leq\epsilon'$, where $\bm{0}$ is the zero matrix and $\epsilon'$ is a distinct hyperparameter governing error (not the same as $\epsilon$). It is readily seen that Eq.~\eqref{eq:oft_principle} attains its global minimum (value zero) if and only if $\bm{W}$ can be obtained from $\bm{W}^0$ by an orthogonal transformation, i.e., $\thickmuskip=2mu \medmuskip=2mu \bm{W} = \bm{R}\bm{W}^0$ with $\bm{R}\in\mathbb{R}^{d\times d}$ an orthogonal matrix (corresponding to rotation for determinant 1 and reflection for determinant $-1$). This reduced step size effectively keeps the updated model close to its pretrained parameters, such that finetuning can be viewed as implicitly minimizing $\thickmuskip=2mu \medmuskip=2mu \| \bm{M} - \bm{M}^0\|$, where $\bm{M}$ represents the finetuned parameters and $\bm{M}^0$ those of the pretrained model. Unlike LoRA, OFT updates neuron weights via layer-wise shared orthogonal transformations in a multiplicative fashion, mathematically guaranteeing the preservation of pairwise angles between neurons within a layer and promoting enhanced stability. The codebase is correct; the results shown in Appendix~\ref{app:rescaled} remain valid.} with $\thickmuskip=2mu \medmuskip=2mu \bm{D}=\text{diag}(s_1,\cdots,s_n)\in\mathbb{R}^{n\times n}$ representing a trainable diagonal matrix whose diagonal entries $s_1,\cdots,s_n$ are constrained to be positive. \textbf{Balancing flexibility and regularization in finetuning}. As illustrated by the randomly selected results in Figure~\ref{fig:control_final}, OFT and COFT consistently generate images that exhibit both high visual fidelity and accurate adherence to control conditions. Additionally, considering that reduced Euclidean distance between the finetuned and pretrained models typically translates to improved retention of pretraining knowledge, we extend this concept with a variant called Constrained Orthogonal Finetuning~(COFT), which restricts the updated model to remain within a fixed-radius hypersphere centered at the pretrained neurons. Additional qualitative results for these three control tasks can be found in Appendix~\ref{app:control_exp}. Alternatively, Stable Diffusion~\cite{rombach2022high} and DALL-E2~\cite{ramesh2022hierarchical} conduct training within a latent space: Stable Diffusion incorporates VQ-GAN~\cite{esser2021taming} to formulate a visual codebook as its latent representation, and DALL-E2 leverages CLIP~\cite{radford2021learning} for constructing a unified latent embedding for both text and images. Notably, the phase spectrum captures angular relationships between input data and basis elements, and it largely encapsulates the core semantic content of the image. Our focus lies on two significant text-to-image finetuning problems: subject-driven generation—where the model generates customized images of a target subject given a handful of subject examples and a prompt—and controllable generation—where the model responds to additional guidance signals. Nonetheless, Cayley parameterization may become parameter-inefficient when dealing with high-dimensional spaces (large $d$), as the orthogonal matrix $\bm{R}$ would then require substantial storage. These findings point to the conclusion that neuron angles (or their directions) are the primary carriers of semantic content in the images. Additional qualitative results are provided in Appendix~\ref{app:more_qualitative}. \end{itemize}

\section{Related Work}

\textbf{Text-to-image diffusion models}. However, textual instructions alone often provide imprecise and limited control over specific visual details in the generated imagery. Furthermore, a theoretical basis for constructing optimal finetuning methods or determining key hyperparameters, such as training duration, is largely absent. \textbf{Finetuning stability and convergence}. The rationale is that adjusting the scale of neurons leaves the hyperspherical energy unchanged, since normalization removes its influence. The OFT paradigm can thus be understood as attempting to minimize the discrepancy in hyperspherical energy between the finetuned and pretrained layers:

\begin{equation}\label{eq:oft_principle}
\footnotesize
\min_{\bm{W}} \left\| \text{HE}(\bm{W})-\text{HE}(\bm{W}^0)\right\|~~\Leftrightarrow~~\min_{\bm{W}} \bigg{\|} \sum_{i\neq j} \|\hat{\bm{w}}_i-\hat{\bm{w}}_j\|^{-1}-\sum_{i\neq j} \|\hat{\bm{w}}^0_i-\hat{\bm{w}}^0_j\|^{-1}\bigg{\|}
\end{equation}

Here, the normalized neuron vectors are defined as $\thickmuskip=2mu \medmuskip=2mu \hat{\bm{w}}_i:=\bm{w}_i/\|\bm{w}_i\|$, and the hyperspherical energy of $\bm{W}$ is given by $\thickmuskip=2mu \medmuskip=2mu \text{HE}(\bm{W}):= \sum_{i\neq j} \|\hat{\bm{w}}_i-\hat{\bm{w}}_j\|^{-1}$. For further experimental validation, see our preliminary work on convolutional OFT finetuning in Appendix~\ref{app:convolution}. Existing methods often operate under the implicit assumption that a smaller Euclidean distance between the parameters of the finetuned and pretrained models correlates with greater retention of the pretrained capabilities. More recently, diffusion models have been introduced to address the same task~\cite{saharia2022palette,voynov2022sketch,wang2022pretraining}. For instance, combining an image’s original phase spectrum with an arbitrary magnitude spectrum enables reconstruction of the image in a way that retains its semantic integrity. Consequently, refining semantic representations in images may be best accomplished by adjusting neuron directions. Distinct from \cite{liu2021orthogonal}, our approach does not seek to reduce hyperspherical energy. 
\begin{abstract}
Large-scale text-to-image diffusion models have demonstrated remarkable prowess in synthesizing photorealistic images conditioned on textual descriptions. OFT is particularly related to LoRA~\cite{hulora2022}, which updates model weights additively under the assumption of a low-rank structure during finetuning. \subsection{Controllable Generation}

\textbf{Settings}. \textbf{Inference efficiency maintained}. Each baseline is tested with optimal hyperparameter choices. To assess controllable generation, we propose a metric for measuring control consistency, which involves extracting the control signal from the output image and comparing it against the original input control. To make this contrast explicit, take a weight matrix sized $\thickmuskip=2mu \medmuskip=2mu 128\times 128$; LoRA, with $\thickmuskip=2mu \medmuskip=2mu r'=8$, entails $2,048$ trainable parameters, whereas OFT (for $\thickmuskip=2mu \medmuskip=2mu r=8$ and no block sharing) requires only $960$. These results underscore the critical role of orthogonality constraints in the OFT method. LoRA, for example, is mainly utilized for Transformer attention weights~\cite{hulora2022}; to ensure a fair comparison, our experiments restrict OFT to these same attention weights. The orthogonality can be maintained using the Cayley transform introduced in Section~\ref{sect:eff_ortho}. \subsection{Efficient Orthogonal Parameterization}\label{sect:eff_ortho}

Conventional approaches for generating orthogonal matrices, such as the differentiable Gram-Schmidt process, can be computationally intensive in practical scenarios~\cite{liu2021orthogonal}. Adopting the protocol of \cite{ruiz2023dreambooth}, we perform a quantitative evaluation measuring subject fidelity (DINO~\cite{caron2021emerging}, CLIP-I~\cite{radford2021learning}), fidelity to text prompts (CLIP-T~\cite{radford2021learning}), and diversity of generated samples (LPIPS~\cite{zhang2018unreasonable}). In the case where $\thickmuskip=2mu \medmuskip=2mu r=1$, the block-diagonal orthogonal structure simplifies to a general unconstrained orthogonal matrix. The standard implementation of OFT displays stable behavior and consistently accurate control by epoch 20. Conversely, reducing $\epsilon$ forces the finetuned model’s behavior to stay closer to that of the pretrained text-to-image diffusion model. The rationale for utilizing orthogonal transformations in neuron finetuning draws, in part, from insights offered by the 2D Fourier transform, which decomposes an image into magnitude and phase spectra. For each approach, we use the model variant with the highest validation CLIP scores. To clarify, we split the motivation into two specific queries: (1) the rationale for adjusting neuron angles (i.e., their directions) during finetuning, and (2) the reasoning behind selecting orthogonal transformations as the mechanism for such angle modifications. Central to this challenge is the absence of a reliable metric for evaluating how well finetuning retains the model's original generative abilities. These observations suggest that OFT achieves superior subject fidelity and controllability simultaneously. \textbf{Qualitative comparison}. At inference, it suffices to premultiply the learned orthogonal matrix $\bm{R}$ with the original pretrained weights $\bm{W}^0$, resulting in the updated weights $\thickmuskip=2mu \medmuskip=2mu \bm{W}=\bm{R}\bm{W}^0$. To address this, we exploit a notable invariance property of hyperspherical energy: applying an identical orthogonal transformation to every neuron in a layer provably maintains their pairwise hyperspherical similarity. \textbf{Model finetuning}. Nonetheless, these methodologies often struggle to maintain the fine identity details of the subject. The main objectives for an effective finetuning technique can be summarized as follows: (1) \emph{training efficiency}, which involves reducing both the number of trainable parameters and the duration of training, and (2) \emph{generalizability preservation}, which seeks to retain the pretrained model’s capacity for producing high-quality and varied outputs. Adapting large pretrained models for downstream tasks has gained traction in recent years~\cite{devlin2018bert,brown2020language,he2022masked}. Distinct from prior finetuning approaches, OFT can formally maintain hyperspherical energy, capturing pairwise neuron relations on the unit hypersphere. Through our experiments, we identify hyperspherical energy as a critical metric for the retention of generative semantics inherent to the pretrained model. Despite their straightforward implementation, these strategies do not guarantee that the generative performance from pretraining will be conserved. \end{abstract}

\section{Introduction}

State-of-the-art text-to-image diffusion models~\cite{saharia2022photorealistic,ramesh2022hierarchical,rombach2022high} have achieved exceptional results in synthesizing high-fidelity images controlled by textual input. \section{Concluding Remarks and Open Problems}

Drawing on the insight that angular relationships between neurons are pivotal in shaping visual semantics, this work introduces orthogonal finetuning—a straightforward yet powerful approach for finetuning text-to-image diffusion models. For illustration, let us consider a fully connected layer $\thickmuskip=2mu \medmuskip=2mu \bm{W}=\{\bm{w}_1,\cdots,\bm{w}_n\}\in\mathbb{R}^{d\times n}$ with individual neuron vectors $\bm{w}_i\in\mathbb{R}^{d}$, and let $\bm{W}^0$ denote its pretrained weights. Within 400 iterations, both DreamBooth and OFT variants successfully provide precise control, while LoRA fails to reliably maintain the subject’s identity. \section{Orthogonal Finetuning}

\subsection{Why Does Orthogonal Transformation Make Sense?}

We begin by motivating the use of orthogonal transformations for finetuning text-to-image diffusion models. OFT, though conceptually straightforward, achieves an effective equilibrium between these two aspects by maintaining pairwise neuron angle relationships. Nevertheless, enforcing the additional $\epsilon$-closeness constraint on the Cayley-transformed matrix is a nontrivial challenge. It remains to be explored whether this collection of orthogonal matrices preserves information from each respective downstream task.